% !TeX program = lualatex
% Compile twice: lualatex 126.tex
% All references and prose are embedded; no BibTeX database or external inputs.
\documentclass[11pt,a4paper]{article}
\usepackage[margin=24mm,headheight=14pt]{geometry}
\usepackage{fontspec}
\setmainfont{Latin Modern Roman}
\setsansfont{Latin Modern Sans}
\setmonofont{Latin Modern Mono}
\usepackage{amsmath,amssymb,mathtools}
\usepackage{unicode-math}
\setmathfont{Latin Modern Math}
\usepackage{microtype}
\usepackage{enumitem,xurl,needspace}
\usepackage[dvipsnames]{xcolor}
\usepackage{hyperref}
\hypersetup{unicode=true,colorlinks=true,linkcolor=MidnightBlue,urlcolor=MidnightBlue,
 pdftitle={500 Mathematical Problem Statements},pdfauthor={Source-annotated compilation},bookmarksnumbered=true}
\usepackage{bookmark}
\usepackage{tocloft}
\setlength{\cftsecnumwidth}{3em}
\cftsetpnumwidth{2.5em}
\cftsetrmarg{4em}
\usepackage{titlesec}
\titleformat{\section}{\Large\bfseries\raggedright}{\thesection}{0.7em}{}
\usepackage{fancyhdr}
\pagestyle{fancy}\fancyhf{}
\fancyhead[L]{\small\sffamily Mathematical problem statements}
\fancyhead[R]{\small Source edition 22}
\fancyfoot[C]{\thepage}
\setlength{\parindent}{0pt}
\setlength{\parskip}{5pt plus 1pt}
\setlength{\emergencystretch}{3em}
\setcounter{tocdepth}{1}
\setcounter{secnumdepth}{1}
\setlist[itemize]{leftmargin=3em,itemsep=5pt,topsep=4pt,parsep=0pt}
\allowdisplaybreaks
\newcommand{\N}{\mathbb{N}}
\newcommand{\Z}{\mathbb{Z}}
\newcommand{\Q}{\mathbb{Q}}
\newcommand{\R}{\mathbb{R}}
\newcommand{\C}{\mathbb{C}}
\newcommand{\F}{\mathbb{F}}
\newcommand{\A}{\mathbb{A}}
\newcommand{\PP}{\mathbb P}
\newcommand{\E}{\mathbb{E}}
\newcommand{\eps}{\varepsilon}
\newcommand{\dd}{\,\mathrm d}
\newcommand{\sref}[2]{\hyperlink{source:#1:#2}{[S#2]}}
\newcommand{\eref}[2]{\hyperlink{extra:#1:#2}{[E#2]}}
% Turn the next switch off for a shorter reading copy. The quotations preserve
% every exactTarget field, including qualifications omitted from a short summary.
\newif\ifshowcatalogue
\showcataloguetrue
\newcommand{\cataloguescope}[1]{\ifshowcatalogue
 \paragraph{Exact scope in the supplied catalogue.}
 \begin{quote}\small #1\end{quote}\fi}
\hypersetup{pdftitle={Problem 126 - Erdos-Hajnal conjecture},pdfauthor={Open Problems Math research notebook}}
\fancyhead[L]{\small\sffamily Open Problems Math \textbar\ Research notebook}
\fancyhead[R]{\small\sffamily 126 \textbar\ 27 September 2026}
\numberwithin{equation}{section}
\begin{document}
\setcounter{section}{125}
% ================================================================
% problemId: problem.erdos-hajnal-conjecture
% source releaseRank: 126; source releaseStatus: open_with_solved_subcases
% exactTarget SHA-256: e68e4f9777c1d2afa90bd0e7fa9cdd05f742889034502483d34946a6eae6ebe9
\clearpage
\section{\texorpdfstring{Erdos-Hajnal Conjecture}{Erdos-Hajnal Conjecture}}\label{126}
\subsection{Definitions and mathematical statement}
An induced copy of $H$ in $G$ is a vertex subset whose inherited edges and nonedges agree with $H$. A clique has all pairwise edges, and an independent set has none. The conjecture is
\[
 \forall H\ \exists c_H>0\ \forall G\text{ with no induced }H:\qquad
 \max(\omega(G),\alpha(G))\ge |V(G)|^{c_H}.
\]
The exponent may depend on the forbidden graph but not on the size of $G$. Excluding ordinary subgraphs or minors is not the same hypothesis.
\par\smallskip\textit{Formulation and concept references:} \sref{126}{1}.
\cataloguescope{Prove that for every fixed finite graph H there is a constant c\_H>0 such that every n-vertex graph with no induced subgraph isomorphic to H contains a clique or an independent set of size at least n\textasciicircum{}\{c\_H\}.}
\subsection{Short English statement}
Does forbidding any fixed induced graph force a polynomially large clique or independent set?
% BEGIN RESEARCH ADDITION 126
\subsection{Research attempt}

\paragraph{Current frontier.}
Buci\'c--Fox--Pham prove the equivalence of the Erd\H{o}s--Hajnal,
polynomial R\"odl, and polynomial Nikiforov properties; the easy
sparsification-to-homogeneous-set direction is particularly useful here.
\eref{126}{1}, Conjecture 2 and Theorem 7.
The August 2026 revision of Huang--Ju--Zhou records the resolution for
all graphs on at most five vertices and develops additional six-vertex
cases.  Its statements do not cover every forbidden finite graph.
\eref{126}{2}, introduction.
The present route uses induced-subgraph exclusion throughout; replacing
it with an ordinary subgraph or minor condition would change the problem.

\paragraph{Attack routes.}
An iterative density argument could succeed with polynomial loss in the
accuracy parameter; ordinary qualitative regularity is not enough.
A recursive decomposition into two large complete or anticomplete
pieces looks stronger and potentially simpler, but it is false even
for a very easy forbidden graph.  A minimal-counterexample argument
would therefore need weaker, quantitatively controlled sparsification,
not pure pairs of fixed linear size.  We prove the quantitative
implication and construct an explicit probabilistic obstruction to the
stronger recursive lemma.

\paragraph{Attempt I: exact exponent from polynomial sparsification.}
Fix $H$.  Suppose there is $a=a_H>0$ such that, for every
$0<\varepsilon<1/2$ and every induced-$H$-free graph on $n$ vertices,
there exists $U$ with
\begin{equation}\label{126:rodl}
 |U|\geq\varepsilon^a n,\qquad
 d(G[U])\leq\varepsilon\quad\text{or}\quad
 d(G[U])\geq1-\varepsilon.
\end{equation}
Density means the fraction of the $\binom{|U|}{2}$ possible edges;
sets with fewer than two vertices can be treated separately.
This is the polynomial R\"odl condition of \eref{126}{1}, not a
hypothesis proved here for every $H$.

For a graph on $m$ vertices with average degree at most $\varepsilon m$,
there is an independent set of size at least
\begin{equation}\label{126:caro}
 \sum_{v}\frac1{\deg(v)+1}\geq\frac{m}{\varepsilon m+1}.
\end{equation}
Indeed, order the vertices uniformly at random and select each vertex
that occurs before all of its neighbors.  The selected vertices are
independent, and the expected count is the displayed sum.  The second
inequality follows from Cauchy--Schwarz.  Applying this to the complement
handles the dense alternative of \eqref{126:rodl}.

Put $t=1/(a+1)$ and choose $\varepsilon=n^{-t}$ once this is below
$1/2$.  Then $m\geq n^t$, and the increasing function
$m\mapsto m/(\varepsilon m+1)$ gives
\[
 \max(\alpha(G),\omega(G))\geq\tfrac12n^t.
\]
To meet the catalogue's coefficient-one normalization, choose a smaller
positive exponent $c<t$, for instance initially $t/2$.  For all large
$n$, $n^c\leq n^t/2$.  Decrease $c$ further if necessary to cover the
finitely many remaining $n\geq2$, using
$\max(\alpha,\omega)\geq2$ there; $n=0,1$ cause no difficulty.
Thus \eqref{126:rodl} implies exactly the supplied conclusion.
The rate and normalization are both explicit; no generic asymptotic
constant has been mistaken for coefficient one.

This calculation also tests slower sparsification.  If the only
available size guarantee were
$|U|\geq n\exp(-C(\log(1/\varepsilon))^2)$, choosing
$\log(1/\varepsilon)$ on the scale $\sqrt{\log n}$ would yield a
homogeneous set of order $\exp(c\sqrt{\log n})$ by the same argument,
not $n^c$.  The gap is a polynomial-versus-superpolynomial loss in the
accuracy parameter.  Merely making a qualitative regularity statement
effective would not close it.

\paragraph{Attempt II: falsify a strong recursive lemma.}
A natural proposed induction asserts that every induced-$H$-free graph
contains disjoint sets of size at least $\delta_H n$ with all or none
of the edges between them.  This is not a valid intermediate lemma
for general $H$, already when $H=K_3$.

Fix $0<\delta<1/2$ and sample $G\sim G(n,p)$ with $p=n^{-3/4}$.
The expected number of triangles is at most $n^{3/4}/6$, so by Markov's
inequality it is at most $n^{7/8}$ with probability tending to one.
There are at most $3^n$ ordered pairs of disjoint vertex sets.  For any
such pair with both sizes at least $\delta n$, the probability of
having no crossing edge is at most
\[
 (1-p)^{\delta^2n^2}\leq\exp(-\delta^2n^{5/4}),
\]
and the probability of having all crossing edges is at most
\[
 p^{\delta^2n^2}=\exp(-\tfrac34\delta^2n^2\log n).
\]
A union bound shows that, with probability tending to one, neither
type of pair exists anywhere in $G$.

Choose a graph with both properties and delete one vertex from each
triangle, deleting at most $n^{7/8}$ vertices in total.  The remaining
graph $G'$ is triangle-free and has $n'=n-o(n)$ vertices.  Any pure
pair in $G'$ is also a pure pair in $G$.  Given any proposed fixed
linear fraction $\eta>0$, run the construction with $\delta=\eta/2$;
then for large $n$, a pair of size $\eta n'$ would exceed $\delta n$
and is impossible.  Hence there is no universal positive pure-pair
fraction for triangle-free graphs.

This does not refute Erd\H{o}s--Hajnal for $K_3$.  Indeed, in a
triangle-free graph every neighborhood is independent.  If the maximum
degree is at least $\sqrt n/2$, this gives a polynomial-size independent
set; otherwise \eqref{126:caro} gives one of order $\sqrt n$.
The counterconstruction rules out the \emph{stronger induction step},
not the much weaker homogeneous-set conclusion.

\paragraph{Attempt III: calibrate where recursive decomposition works.}
For graphs built from one-vertex graphs by disjoint unions and complete
joins, one can prove the useful invariant
\begin{equation}\label{126:product}
 \alpha(G)\omega(G)\geq|V(G)|.
\end{equation}
For a disjoint union, $\alpha$ is the sum of the component independence
numbers and $\omega$ their maximum, so the product is at least the
sum of their products.  For a complete join, interchange these roles.
Induction proves \eqref{126:product}, hence a homogeneous set of size
at least $\sqrt n$.  This is a verified decomposition-based subcase.
It does not assert that every $H$-free graph admits such a decomposition.
The probabilistic example explains why trying to impose an equally
strong recursive structure on all forbidden-induced-graph classes is
not justified.

A more promising next step would be a density-reduction process whose
logarithmic size loss is $O_H(\log(1/\varepsilon))$ in total, rather
than a loss incurred anew at each accuracy scale.  Such an amortized
bound would imply \eqref{126:rodl}.  No process with that bound for
arbitrary $H$ is proved here; the known equivalence in \eref{126}{1}
warns that proving it is essentially the core problem, not a routine
regularity estimate.

\paragraph{Stress test.}
The random-graph alteration destroys every triangle while preserving
the absence of large pure pairs, since deleting vertices cannot create
one.  The relevant probability estimates dominate the $3^n$ choices;
no heuristic independence among different pairs is assumed.  The
sparsification calculation uses the complement only to extract a
clique, not to change which graph $H$ is forbidden.  Constants and small
$n$ are handled explicitly.

The companion script checks \eqref{126:caro} and
\eqref{126:product} on small exact graphs.  A finite random sample is
not used to certify the asymptotic counterexample: the union and Markov
bounds above supply that proof.  No stronger claim of a counterexample
to Erd\H{o}s--Hajnal follows from these method tests.

\paragraph{Outcome.}
\textbf{Outcome: Conditional reduction.}

The argument quantifies the sparsification condition sufficient for
the exact target and proves why a tempting linear-size pure-pair
substitute fails.  It also verifies a recursive class where the product
invariant succeeds.  Polynomial sparsification for every forbidden
$H$ remains unproved here, so no extension of the known forbidden-graph
range is claimed.  These deductions are not presented as historically
new results.
% END RESEARCH ADDITION 126
\subsection{Sources}
\begin{itemize}\raggedright
\item[\textnormal{[S1]}]\hypertarget{source:126:1}{} Lior Gishboliner and Ethan Honest, Disperse hypergraphs, Introduction (2025).
\url{https://www.cambridge.org/core/journals/combinatorics-probability-and-computing/article/disperse-hypergraphs/49C68DA657349D5D7BAB81EE3E971729}.
{\small\textit{Catalogue formulation source.}}
% BEGIN RESEARCH REFERENCES 126
\item[\textnormal{[E1]}]\hypertarget{extra:126:1}{} Matija Buci\'c, Jacob Fox, and Huy Tuan Pham.
\emph{Equivalence between Erd\H{o}s--Hajnal and polynomial R\"odl and Nikiforov conjectures}, arXiv:2403.08303v2, 19 April 2024. Conjecture 2 and Theorem 7; the quantitative sparsification formulation and equivalence.
\url{https://arxiv.org/html/2403.08303v2}.
{\small\textit{Consulted 27 September 2026 for the indicated result or scope.}}
\item[\textnormal{[E2]}]\hypertarget{extra:126:2}{} Shenwei Huang, Yiao Ju, and Yidong Zhou.
\emph{Erd\H{o}s--Hajnal conjecture beyond five-vertex graphs}, arXiv:2606.06258v3, 31 August 2026. Introduction and stated additional six-vertex cases. Recent preprint; not independently proof-audited here.
\url{https://arxiv.org/html/2606.06258v3}.
{\small\textit{Consulted 27 September 2026 for the indicated result or scope.}}
% END RESEARCH REFERENCES 126
\end{itemize}


\end{document}
